\documentclass[prl,twocolumn,nofootinbib]{revtex4}
\usepackage{bm,amsmath,amssymb,graphicx}
%\usepackage{pgfplots}
%\pgfplotsset{width=10cm,compat=1.9}
%\documentclass{article}
%\usepackage[utf8]{inputenc}

%\title{Jordan and Niko research}
%\author{npoplawski }
%\date{March 2019}

\begin{document}

\title{Black Hole Genesis and origin of inertia}

\author{Nikodem Pop{\l}awski}
\altaffiliation{NPoplawski@newhaven.edu}

\affiliation{Department of Mathematics and Physics, University of New Haven, 300 Boston Post Road, West Haven, CT 06516, USA}


\begin{abstract}
I propose that if the universe was born as a baby universe on the other side of the event horizon of a black hole existing in a parent universe, then the corresponding white hole provides the absolute inertial frame of reference in the universe.
The principle of relativity then allows to construct an infinity of other inertial frames.
Consequently, this scenario could give the origin of inertia and complete Einstein's general theory of relativity by making it consistent with Mach's principle.
\end{abstract}
\maketitle


If there were only one body in the universe, for example Earth, then it would have no relative motion with respect to other bodies.
Yet, a Foucault's pendulum could determine Earth's rotation about its own axis \cite{LL1}.
But, with respect to what would Earth rotate if there were no other bodies?
Earth's state of motion would have no meaning in that case.

Newton's rotating bucket argument demonstrated that true rotational motion cannot be defined as the relative rotation of the body with respect to the immediately surrounding bodies \cite{Newton}.
More generally, true motion and rest cannot be defined relative to other bodies. Instead, they can be defined only by reference to absolute space.

Einstein's general theory of relativity, in which the motion of bodies is determined by the local geometry of spacetime \cite{LL2}, reduces absolute space and time to local geodesics that are sufficient to describe this geometry.
Absolute space becomes a field that is described by the metric tensor.
True motion and rest are defined by reference to the metric tensor that asymptotically (far away from physical bodies) tends to the form determined by the condition of constant curvature, which depends on whether the universe is flat, closed, or open.

According to Einstein, the metric tensor is determined locally by the distribution of matter.
What determines the asymptotic form of the metric tensor that took the role of absolute space?
Mach's principle states that the overall distribution of matter provides absolute space: local physical laws are determined by the large-scale structure of the universe \cite{Mach}.
Consequently, the motion of the distant stars determines the local inertial frame.
But, if there were no bodies other than Earth, there would be no distant matter that could determine the metric tensor.
The metric would be redundant because it would have nothing to relate to.
Yet, the metric, taking the role of absolute space, must exist in order to explain the difference between two scenarios in which the plane of oscillation of a Foucault’s pendulum rotates (indicating Earth's rotation with respect to the metric field) or not.

Therefore, there must exist a body in the universe that determines absolute space.
I propose that Black Hole Genesis (BHG) provides a natural answer \cite{ApJ}.
If our universe was born as a closed, baby universe \cite{BHG} formed on the other side of the event horizon of a parent black hole existing in a parent universe, then that black hole would be seen in the baby universe as a primordial white hole.
That white hole determines absolute space: the frame of reference in which the white hole is at rest is the absolute inertial frame of reference (AIFR).\footnote{
The white hole being at rest means that its comoving distance (fixed relative to growing space) from the observer does not change.
The physical distance will grow because the universe is expanding.
}
This frame defines the absolute time, called the cosmic time, which appears in the Friedmann equations of cosmology.
It also defines absolute simultaneity and comoving distances.

Newton's first law of dynamics (the law of inertia) has two parts: 1) There exist inertial frames of reference, and 2) In an inertial frame, a free body (without forces acting on it) has a constant velocity: “a body at rest stays at rest and a body in motion stays in motion.”
In the Lagrangian formulation, an inertial frame of reference is a frame in which space is homogeneous and isotropic, and time is homogeneous \cite{LL1}.
I propose that the primordial white hole guarantees that at least one inertial frame exists: the AIFR, thus explaining the origin of inertia.
The principle of relativity of Galileo and Einstein then allows to construct, through the Lorentz transformations, an infinity of other inertial frames that are mechanically equivalent to AIFR and to one another \cite{LL2}.

I propose the following conjecture: the AIFR is also the frame in which the total momentum and total angular momentum of the matter in the universe are zero (they can be measured only within the observer's cosmological horizon).
These definitions are determinate because the universe formed by a black hole is closed (with the exception of the white hole that connects the universe to the parent universe through an Einstein--Rosen bridge).
This conjecture has the spirit of Mach's principle: all distant matter determines inertia \cite{Mach}.
I also propose another conjecture: the AIFR coincides with the frame in which the cosmic microwave background radiation is isotropic (on average, without accounting for tiny fluctuations that have led to large-structure formation) \cite{pref}.

Therefore, BHG completes Einstein's general theory of relativity by making it Machian.
The parent black hole, seen in our universe as a white hole, constitutes the Machian distant matter that determines absolute space and provides inertial frames of references.

There are further implications of BHG.
If our universe was born on the other side of the event horizon of a black hole existing in a parent universe, then every black hole creates a new, baby universe.
These universes form a multiverse.
However, they are not parallel.
An object can exist at any moment of its timeline (measured in its proper time defined in its rest frame) only in one universe.

Since the motion of matter through an event horizon can only occur in one direction, that motion can define the past and future.
This existence of the arrow of time at the event horizon can be continuously extended to all other points in space.
It will also be extended to the cosmic time, defined as the time in AIFR, which is a coordinate that measures the expansion of the universe.
Accordingly, the cosmic arrow of time in the universe would be inherited from the parent universe in BHG \cite{essay}.

The second law of thermodynamics states that the total entropy of an isolated system can never decrease over time.
Once the total entropy of a system and its surroundings reaches a maximum, it would remain constant: the system is in thermodynamic equilibrium with the surroundings.
A black hole, which forms in the infinite future as measured in AIFR, is a cosmic example of such a system.
However, AIFR of the parent universe cannot be extended beyond the event horizon of a black hole because of the infinite redshift at the horizon \cite{LL2}.
On the other side of the event horizon, a baby universe has its own AIFR and its own cosmic time.
In that growing universe, entropy can increase further \cite{ApJ}.
Therefore, the thermodynamic and cosmic arrows of time coincide.

The black hole information paradox does not exist in this scenario.
The information goes from the parent universe to a baby universe formed on the other side of the black hole's event horizon.
Since the curvature of the closed universe is absolute, the gravitational force is geometrical and does not need a mediating particle: the graviton does not exist.

A physical law that turns black holes into Einstein--Rosen bridges to new, baby universes must avoid the black-hole singularity.
The simplest and most natural mechanism for preventing gravitational singularities is provided by spacetime torsion within the Einstein--Cartan theory of gravity \cite{Newton,EC}.
In this theory, torsion is coupled to the intrinsic angular momentum of fermionic matter, allowing for the spin-orbit interaction that follows from the Dirac equation.
Accordingly, torsion brings the consistency between relativistic quantum mechanics and curved spacetime.
At extremely high densities, torsion manifests itself as repulsive gravity, preventing the formation of a singularity and creating a Big Bounce that starts a new universe \cite{ApJ,avert}.

The inertia in the universe may therefore originate from the universe being formed by a black hole that naturally provides AIFR: the absolute inertial frame of reference.
That formation is physically realized through initial gravitational attraction from curvature, which is later countered by gravitational repulsion from torsion.


%\section*{Acknowledgments}
%{\it Acknowledgment.}
This work was funded by the University Research Scholar program at the University of New Haven.

\begin{thebibliography}{}
\bibitem{LL1} L. D. Landau and E. M. Lifshitz, {\it Mechanics} (Pergamon, 1976).
\bibitem{Newton} M. Born, {\em Einstein's Theory of Relativity} (Dover, 1962).
\bibitem{LL2} E. Schr\"{o}dinger, {\it Space-time Structure} (Cambridge University Press, 1954); L. D. Landau and E. M. Lifshitz, {\it The Classical Theory of Fields} (Pergamon, 1975).
\bibitem{Mach} S. Weinberg, {\em Gravitation and Cosmology} (Wiley, 1972); E. A. Lord, {\em Tensors, Relativity and Cosmology} (McGraw-Hill, 1976); I. Ciufolini and J. A. Wheeler, {\em Gravitation and Inertia} (Princeton University Press, 1995).
\bibitem{ApJ} N. J. Pop{\l}awski, Phys. Lett. B {\bf 694}, 181 (2010); Phys. Lett. B {\bf 701}, 672 (2011); N. Pop{\l}awski, Astrophys. J. {\bf 832}, 96 (2016).
\bibitem{BHG} I. D. Novikov, J. Exp. Theor. Phys. Lett. {\bf 3}, 142 (1966); R. K. Pathria, Nature {\bf 240}, 298 (1972); V. P. Frolov, M. A. Markov, and V. F. Mukhanov, Phys. Lett. B {\bf 216}, 272 (1989); Phys. Rev. D {\bf 41}, 383 (1990); L. Smolin, Class. Quantum Grav. {\bf 9}, 173 (1992); S. Hawking, {\it Black Holes and Baby Universes and other Essays} (Bantam Dell, 1993); W. M. Stuckey, Am. J. Phys. {\bf 62}, 788 (1994); D. A. Easson and R. H. Brandenberger, J. High Energ. Phys. {\bf 06}, 024 (2001); J. Smoller and B. Temple, Proc. Natl. Acad. Sci. USA {\bf 100}, 11216 (2003); N. J. Pop{\l}awski, Phys. Lett. B {\bf 687}, 110 (2010).
\bibitem{pref} P. J. E. Peebles, {\em Principles of Physical Cosmology} (Princeton University Press, 1993).
\bibitem{essay} N. Pop{\l}awski, Int. J. Mod. Phys. D {\bf 27}, 1847020 (2018).
\bibitem{EC} T. W. B. Kibble, J. Math. Phys. {\bf 2}, 212 (1961); D. W. Sciama, in {\em Recent Developments in General Relativity}, p. 415 (Pergamon, 1962); Rev. Mod. Phys. {\bf 36}, 463 (1964); Rev. Mod. Phys. {\bf 36}, 1103 (1964); F. W. Hehl, P. von der Heyde, G. D. Kerlick, and J. M. Nester, Rev. Mod. Phys. {\bf 48}, 393 (1976); V. de Sabbata and M. Gasperini, {\it Introduction to Gravitation} (World Scientific, 1985); K. Nomura, T. Shirafuji, and K. Hayashi, Prog. Theor. Phys. {\bf 86}, 1239 (1991); V. de Sabbata and C. Sivaram, {\it Spin and Torsion in Gravitation} (World Scientific, 1994); N. J. Pop{\l}awski, arXiv:0911.0334; Phys. Lett. B {\bf 690}, 73 (2010); Phys. Lett. B {\bf 727}, 575 (2013).
\bibitem{avert} W. Kopczy\'{n}ski, Phys. Lett. A {\bf 39}, 219 (1972); Phys. Lett. A {\bf 43}, 63 (1973); A. Trautman, Nature (Phys. Sci.) {\bf 242}, 7 (1973); F. W. Hehl, P. von der Heyde, and G. D. Kerlick, Phys. Rev. D {\bf 10}, 1066 (1974); B. Kuchowicz, Gen. Relativ. Gravit. {\bf 9}, 511 (1978); M. Gasperini, Phys. Rev. Lett. {\bf 56}, 2873 (1986); Y. N. Obukhov and V. A. Korotky, Class. Quantum Grav. {\bf 4}, 1633 (1987); Gen. Relativ. Gravit. {\bf 44}, 1007 (2012); N. Pop{\l}awski, Phys. Rev. D {\bf 85}, 107502 (2012); S. Desai and N. J. Pop{\l}awski, Phys. Lett. B {\bf 755}, 183 (2016); G. Unger and N. Pop{\l}awski, Astrophys. J. {\bf 870}, 78 (2019); J. L. Cubero and N. J. Pop{\l}awski, arXiv:1906.11824.
\end{thebibliography}

\end{document}
